\chapter{Curved transition layers and the Jacobi obstruction}\label{sec:phase-curved-transition}
We follow the local calculations of \cite{ChAC}, Appendix B, Problem 3, pp.~29--31, with explicit signs, boundedness conditions and error norms. Let $\Sigma^2\subset\mathbb R^3$ be a smooth two-sided embedded surface, $\nu$ a chosen unit normal, and $A=D\nu$. Work on a relatively compact coordinate patch $K$ inside a larger patch $U$ on which
\[
 F(y,z)=y+z\nu(y),\qquad |z|<\rho,
\]
is an injective tubular parametrization, with $I+zA$ uniformly positive definite on a neighborhood of $K\times[-\rho,\rho]$. All estimates below are uniform on this fixed tube over $K$, with constants depending on its geometry and the stated smooth functions. We write $H=\operatorname{tr}A$, $K_\Sigma=\det A$, and retain the quartic potential, profile $q$, and $\psi=q'$ from Section \ref{sec:phase-linearized-transition}. No closed nonflat minimal surface in Euclidean space is assumed.

\begin{lemma}[The exact metric, Jacobian and parallel curvature]\label{lem:phase-exact-fermi-geometry}
In these coordinates the Euclidean metric, volume density relative to $d\mu_\Sigma\,dz$, and parallel shape operator are
\[
 g_{\mathbb R^3}=dz^2+g_z,\qquad
 g_z(v,w)=g((I+zA)v,(I+zA)w),\qquad
 D(y,z)=\det(I+zA)=1+Hz+K_\Sigma z^2,
\]
\[
 A_z=A(I+zA)^{-1},\qquad
 H_z=\partial_z\log D=\sum_{i=1}^2\frac{\kappa_i}{1+z\kappa_i},
 \qquad \partial_zH_z=-|A_z|^2.
\]
Here $\kappa_1,\kappa_2$ are the principal curvatures at $y$, and $A_z$ is identified through the parallel tangent planes.
\end{lemma}
\begin{proof}
Differentiate $F$ tangentially to get $dF(v)=(I+zA)v$, while $\partial_zF=\nu$ is orthogonal to every tangential derivative and has length one. This proves the metric formula. Diagonalizing $A$ gives the positive volume ratio $(1+z\kappa_1)(1+z\kappa_2)$ and its stated polynomial. The normal of each parallel surface is still $\nu(y)$, so differentiation along its tangent vector $(I+zA)v$ gives $Av$. Thus its shape operator is $A(I+zA)^{-1}$. Taking its trace and differentiating each eigenvalue proves both remaining formulas. Positivity and injectivity of the chosen tube are used to avoid a focal point or a self-intersection, not inferred merely from the formulas.
\end{proof}

\begin{proposition}[The Laplacian and the special cancellation for minimal surfaces]\label{prop:phase-fermi-laplacian}
For every smooth function $f$ in the tube,
\[
 \Delta f=f_{zz}+H_zf_z+\Delta_{g_z}f.
\]
If $\Sigma$ is minimal and $a(y)=|A(y)|^2$, then
\[
 D=1-\frac{a z^2}{2},\qquad
 H_z=-\frac{az}{1-az^2/2}=-az+O(z^3).
\]
For a smooth function $h(y)$ independent of $z$,
$\Delta_{g_z}h=\Delta_\Sigma h+O(|z|)$ uniformly on the fixed patch, with the last constant controlled by its first two derivatives and the local geometry.
\end{proposition}
\begin{proof}
In block diagonal coordinates the divergence formula for the Laplacian separates into
$D^{-1}\partial_z(Df_z)$ and the tangential Laplace--Beltrami operator. Lemma \ref{lem:phase-exact-fermi-geometry} gives the coefficient of $f_z$. For a minimal surface the principal curvatures are $\kappa,-\kappa$, hence $a=2\kappa^2$ and $D=1-az^2/2$. Differentiating its logarithm gives the exact rational expression. Its difference from $-az$ is $-a^2z^3/[2(1-az^2/2)]$, uniformly of order $z^3$ on the chosen tube. Finally the coefficients of $\Delta_{g_z}$ and their first tangential derivatives vary smoothly with $z$. Subtract their values at zero in the coordinate formula; every difference is bounded by $C|z|$. This gives the last estimate without asserting an equality between the metrics at nonzero $z$.
\end{proof}

\begin{proposition}[When the distance profile is an exact solution]\label{prop:phase-distance-profile-flatness}
Put $u_\varepsilon(F(y,z))=q(z/\varepsilon)$ and
$R_\varepsilon=\varepsilon^2\Delta u_\varepsilon-W'(u_\varepsilon)$.
Then $R_\varepsilon=\varepsilon H_z\psi(z/\varepsilon)$, which is uniformly $O(\varepsilon)$, and is uniformly $O(\varepsilon^2)$ when $\Sigma$ is minimal. If this profile is an exact solution throughout a tube for even one $\varepsilon>0$, then $A=0$ on its base; every connected base patch is contained in a plane. If $H(y_0)\ne0$, the general $O(\varepsilon)$ estimate is sharp at $(y_0,0)$.
\end{proposition}
\begin{proof}
The equation $q''=W'(q)$ cancels the normal second derivative, leaving precisely the residual already identified in Proposition \ref{prop:phase-recovery-residual}. The boundedness of $H_z$ gives the general estimate. In the minimal case $|H_z|\leq C|z|$, and $|t|\psi(t)$ is bounded, giving $C\varepsilon^2$. If the residual vanishes identically, positivity of $\psi$ forces $H_z=0$ for every $z$ in the tube. Differentiate at zero using Lemma \ref{lem:phase-exact-fermi-geometry}: $0=-|A|^2$. The normal is therefore constant on each connected patch, and $y\cdot\nu$ is constant along it, proving planarity. Finally $R_\varepsilon(y_0,0)=\varepsilon H(y_0)\psi(0)$ proves sharpness when the mean curvature does not vanish.
\end{proof}

\begin{lemma}[The bounded odd normal corrector]\label{lem:phase-normal-curvature-corrector}
There is exactly one bounded solution $j$ of
\[
 j''-Vj=t\psi,\qquad j(0)=0.
\]
It is odd, belongs to $H^2(\mathbb R)$, and is given by
\[
 j(t)=\psi(t)\int_0^t\psi(s)^{-2}
       \left(\int_{-\infty}^s\tau\psi(\tau)^2\,d\tau\right)ds.
\]
It is negative for $t>0$, orthogonal to $\psi$, and satisfies
$|j|+|j'|+|j''|\leq C(1+t^2)e^{-\sqrt2|t|}$.
\end{lemma}
\begin{proof}
The total integral of $\tau\psi(\tau)^2$ is zero by oddness. Its primitive $I(s)$ from minus infinity is even and strictly negative at every finite $s$: for $s\leq0$ this is immediate, and for $s>0$ use $I(s)=-\int_s^\infty\tau\psi(\tau)^2d\tau$. Thus $\psi^{-2}I$ is even and negative, proving the claimed parity and sign of the formula. Writing $j=\psi v$ gives
\[
 (\psi^2v')'=t\psi^2,\qquad j''-Vj=\psi v''+2\psi'v'=t\psi,
\]
so the formula solves the correct equation with the prescribed value at zero.

For $s\geq0$, the exponential bounds on $\psi$ give
$|I(s)|\leq C(1+s)e^{-2\sqrt2s}$ and $\psi(s)^{-2}|I(s)|\leq C(1+s)$. Therefore $|v(t)|\leq C(1+t^2)$ and $|v'(t)|\leq C(1+|t|)$; parity handles the other end. These bounds give the estimates for $j$ and $j'$, and the differential equation gives the bound for $j''$. They imply $j\in H^2$ and justify its orthogonality to the even function $\psi$. The difference of two bounded solutions is a multiple of $\psi$ by Proposition \ref{prop:phase-bounded-normal-kernel}; its zero value at the origin forces that multiple to vanish. Both boundedness and the value at zero are needed for this uniqueness statement. The linear coefficient is $V=W''(q)$, as dictated by linearization of $W'$.
\end{proof}

\begin{proposition}[A fourth-order residual for the unshifted corrected layer]\label{prop:phase-corrected-layer-residual}
Assume $\Sigma$ is minimal and set $a=|A|^2$, $t=z/\varepsilon$ and
\[
 u_\varepsilon=q(t)+\varepsilon^2a(y)j(t).
\]
Then $\|\varepsilon^2\Delta u_\varepsilon-W'(u_\varepsilon)\|_{C^0(K\times(-\rho,\rho))}
\leq C\varepsilon^4$ for $0<\varepsilon\leq1$.
This is a local approximate solution, without imposed boundary values or a claim that it can already be corrected to an exact solution.
\end{proposition}
\begin{proof}
Since $W'(u)=u^3-u$, direct substitution and Lemma \ref{lem:phase-normal-curvature-corrector} give the exact residual
\[
 \varepsilon(H_z+az)\psi(t)+
 \varepsilon^3H_z a j'(t)+\varepsilon^4j(t)\Delta_{g_z}a
 -3\varepsilon^4q(t)a^2j(t)^2-\varepsilon^6a^3j(t)^3.
\]
The first term is bounded by $C\varepsilon|z|^3\psi(t)\leq C\varepsilon^4$, because $|t|^3\psi(t)$ is bounded. In the second term use $|H_z|\leq C|z|$ and boundedness of $tj'$. The remaining terms are of order at least four by bounded geometry and the corrector estimates. This proves the asserted uniform zeroth-order bound. The improved order uses the odd expansion of $H_z$ for a two-dimensional minimal surface in Euclidean three-space; it is not an asserted dimension-independent cancellation.
\end{proof}

\begin{lemma}[Shifted layers and the cubic error]\label{lem:phase-shifted-layer-error}
On a minimal patch let $h$ be smooth and put
\[
 u_\varepsilon(y,z)=q(t-\varepsilon h(y))+\varepsilon^2a(y)j(t).
\]
With $a=|A|^2$ and $t=z/\varepsilon$, its residual is
\[
 R_\varepsilon(y,\varepsilon t)=\varepsilon^3
 \left[-\psi(t)\Delta_\Sigma h+
 ah\{t\psi'(t)+W'''(q(t))\psi(t)j(t)\}\right]+O(\varepsilon^4)
\]
uniformly for $y\in K$, $|t|<\rho/\varepsilon$. The displacement of the profile's center is $\varepsilon^2h$, not $\varepsilon h$.
\end{lemma}
\begin{proof}
Write $s=t-\varepsilon h$. Tangential differentiation gives
$\Delta_{g_z}q(s)=-\varepsilon\psi(s)\Delta_{g_z}h+\varepsilon^2q''(s)|D_{g_z}h|^2$.
Normal differentiation and expansion of the cubic potential give
\[
 R_\varepsilon=\varepsilon H_z\psi(s)
 -\varepsilon^3\psi(s)\Delta_{g_z}h
 +\varepsilon^4q''(s)|D_{g_z}h|^2
 +\varepsilon^2a\{j''-W''(q(s))j\}
 +\varepsilon^3H_zaj'+\varepsilon^4j\Delta_{g_z}a
 -3\varepsilon^4q(s)a^2j^2-\varepsilon^6a^3j^3.
\]
Use $H_z=-a\varepsilon t+O(\varepsilon^3|t|^3)$,
$\psi(s)=\psi(t)-\varepsilon h\psi'(t)+O(\varepsilon^2)$, and
$W''(q(s))=V(t)-\varepsilon hW'''(q(t))\psi(t)+O(\varepsilon^2)$.
The two terms $-\varepsilon^2at\psi$ and $+\varepsilon^2at\psi$ cancel. The remaining cubic terms are precisely the displayed expression. The difference $\Delta_{g_z}h-\Delta_\Sigma h$ contributes $O(\varepsilon|t|)$ multiplied by $\varepsilon^3\psi(s)$, hence only $O(\varepsilon^4)$.

Uniformity on the expanding interval in $t$ follows from exponential decay of every positive-order derivative of $q$ and the polynomial-times-exponential bounds on $j$. Taylor remainders for a shift of size at most $\varepsilon\|h\|_\infty$ have the same decay with a fixed larger constant; factors $|t|^m$ are therefore harmless. These are $C^0$ estimates in the physical tube. Unscaled normal derivatives introduce powers of $\varepsilon^{-1}$, so no assertion about errors of the same order in every unscaled derivative is made. Finally the center equation $s=0$ gives $z=\varepsilon^2h$.
\end{proof}

\begin{proposition}[Projection onto translation produces the surface Jacobi operator]\label{prop:phase-jacobi-obstruction}
Under Lemma \ref{lem:phase-shifted-layer-error},
\[
 \int_{-\rho/\varepsilon}^{\rho/\varepsilon}
 R_\varepsilon(y,\varepsilon t)\psi(t)\,dt
 =-\varepsilon^3c_W(\Delta_\Sigma h+|A|^2h)+O(\varepsilon^4)
\]
uniformly for $y\in K$. Thus cancellation of the leading projected residual is the Jacobi equation on the interface. This is the solvability calculation in \cite{ChAC}, Problem 3(o)--(p); it does not supply a global inverse or a nonlinear gluing theorem.
\end{proposition}
\begin{proof}
All endpoint terms below vanish by the corrector and profile bounds. First,
\[
 I_1=\int t\psi'\psi=-\frac12\int\psi^2=-\frac{c_W}{2}.
\]
Differentiating $\psi''-V\psi=0$ gives
$\psi'''-V\psi'=V'\psi=W'''(q)\psi^2$.
The self-adjoint one-dimensional integration by parts then yields
\[
 I_2=\int W'''(q)\psi^2j
 =\int(\psi'''-V\psi')j
 =\int\psi'(j''-Vj)=\int t\psi'\psi=-\frac{c_W}{2}.
\]
Pair the cubic residual with $\psi$: its coefficient is
$-c_W\Delta_\Sigma h+ah(I_1+I_2)=-c_W(\Delta_\Sigma h+ah)$.
Replacing the finite interval by the line in this coefficient incurs only an exponentially small error. The uniform remainder, paired with the integrable $\psi$, is still $O(\varepsilon^4)$, without a factor equal to the length of the interval. This proves the formula and identifies exactly which geometric operator arises.
\end{proof}
